\answer{ex:09:1}{(a)~$P_\infty=0.2$、記憶$20\,\text{s}$。(b)~$P_\infty=1$、記憶$1\,\text{s}$。(c)~記憶は雑音の振幅に比例するので、6倍。}
\answer{ex:09:2}{$P(t)=P_\infty\coth(t\sqrt{q/R})$。(a)~$\frac12\sqrt{R/q}=5\,\text{s}$で、記憶の半分。(b)~$t=\frac12\ln21\,\sqrt{R/q}\approx15\,\text{s}$。高い\nt{ACET}はこれだけ保たれなければならない。}
\answer{ex:09:3}{分散は$qT/2+R/(2T)$。$T^*=\sqrt{R/q}$で分散は$\sqrt{qR}=P_\infty$、フィルタ~\eqref{eq:kalman}と同じ。増加は$(\kappa+1/\kappa)/2$倍で、$1.25$と$5.05$。}
\answer{ex:09:4}{$a>a_w=1-\theta/(mw)$：$w=0.041$で$0.15$、$w=0.045$で$0.22$。}
\answer{ex:09:5}{誤りは$M\approx40$--$60$で現れ、$M=80$で余分なニューロン$1.9$個、$M=100$でパターンの割合$0.86$。他のパターンからの干渉の分散は$\approx(M-1)p(1-p)/N$なので、$M_{\max}\approx80$、$M_{\max}\propto N/p$。}
\answer{ex:09:6}{たとえば$\eta(a)=\eta\min\bigl(1,\max(0,(a-0.3)/0.6)\bigr)$。$\eta a$では無関係なニューロン3個がすべてパターンに入った。$\eta(a)$では1個も入らず、$a\ge0.9$での書き込みは同じ（割合$1.00$）。}
\answer{ex:09:7}{(a)~$40$回。問題~\ref{ex:09:6}の$\eta(a)$では決して蓄積しない。(b)~$200$回の再生で、$10$晩。}
